% Progress Report I — due Nov 17, 2 pages, 11 pt.
% Build: make report1
\documentclass[11pt]{article}
\usepackage{deqreport}
\reporttitle{Progress Report I}

\begin{document}
\makeheader

\section{Summary}

The hand-derived NumPy implementation of a fully connected deep equilibrium
model (DEQ) is complete. Its gradients match central differences to a worst
relative error of $1.3 \times 10^{-6}$, confirming that the implicit backward
pass is correct, and the model reaches 97.5\% test accuracy on MNIST after
five epochs. Both planned forward solvers are implemented. The project is on
schedule for the convolutional PyTorch model in Progress Report~II.

\section{What Was Built}

\subsection{Model and forward solver}

The equilibrium layer is $\fth(z, x) = \ReLU(zW^{\!\top} + xU^{\!\top} + b)$,
and the forward pass solves
\begin{equation}
  \zstar = \fth(\zstar, x)
  \label{eq:fp}
\end{equation}
with two solvers written from scratch: plain fixed-point iteration, and
Anderson acceleration~\cite{anderson1965iterative} with memory $m = 5$, which
keeps the last $m$ iterates and takes the step minimising the residual over
their affine hull. Both stop at a relative residual below $10^{-3}$ or after
50 iterations. At initialisation $W$ is rescaled to spectral norm $0.9$ so the
layer begins as a contraction; without this the solver stalls on the first
batch.

\subsection{Backward pass derived by hand}

Differentiating Eq.~\eqref{eq:fp} implicitly gives
\begin{equation}
  \dd{\ell}{\theta}
  = \dd{\ell}{\zstar}\left(I - \dd{\fth}{\zstar}\right)^{-1}\dd{\fth}{\theta}.
  \label{eq:implicit}
\end{equation}
Rather than forming the inverse, the adjoint $u = \dd{\ell}{\zstar}(I-J)^{-1}$
is obtained as the fixed point of $u = uJ + \dd{\ell}{\zstar}$, which is solved
with the same two solvers. Writing $a$ for the pre-activation and
$M = \mathbf{1}[a > 0]$, the Jacobian is $\partial f_i / \partial z_j = M_i
W_{ij}$, so the two products the backward pass needs are
\begin{equation}
  uJ = (u \odot M)\,W,
  \qquad
  \dd{\ell}{W} = (u \odot M)^{\!\top} \zstar,
  \qquad
  \dd{\ell}{U} = (u \odot M)^{\!\top} x .
\end{equation}
These are derived and coded by hand; no automatic differentiation is used
anywhere in the implementation.

\section{Results So Far}

\begin{table}[h]
  \centering
  \caption{Central-difference gradient check, $\varepsilon = 10^{-6}$, solver
    tolerance $10^{-10}$. Relative error is taken over entries where
    $\max(|g_{\text{analytic}}|, |g_{\text{numeric}}|) > 10^{-8}$; with a ReLU
    layer most entries have a true gradient of order $10^{-10}$ and are judged
    on absolute error alone.}
  \begin{tabular}{@{}lrrr@{}}
    \toprule
    Parameter & Max abs.\ error & Max rel.\ error & Entries checked \\
    \midrule
    $W$ & $6.5\times10^{-11}$ & $1.27\times10^{-6}$ & 9/25 \\
    $U$ & $1.1\times10^{-10}$ & $1.35\times10^{-8}$ & 30/30 \\
    $b$ & $3.4\times10^{-11}$ & $1.86\times10^{-9}$ & 5/5 \\
    $V$ & $6.9\times10^{-11}$ & $3.29\times10^{-7}$ & 15/15 \\
    $c$ & $9.4\times10^{-11}$ & $1.53\times10^{-9}$ & 3/3 \\
    \bottomrule
  \end{tabular}
  \label{tab:gradcheck}
\end{table}

\begin{table}[h]
  \centering
  \caption{MNIST, 5 epochs, hidden width 128, Adam at $10^{-3}$, batch 128.
    Iteration counts are means over the final epoch.}
  \begin{tabular}{@{}lrrrr@{}}
    \toprule
    Forward solver & Test acc.\ (\%) & Fwd.\ iters & Bwd.\ iters & Time/epoch (s) \\
    \midrule
    Fixed-point        & 97.54 & 17.0 & 17.6 & 5.1 \\
    Anderson ($m = 5$) & 97.18 & 12.9 & 11.7 & 10.1 \\
    \bottomrule
  \end{tabular}
  \label{tab:mnist}
\end{table}

Table~\ref{tab:gradcheck} confirms the implicit gradient: every parameter
agrees with central differences to at least six significant figures on the
entries that carry signal. Table~\ref{tab:mnist} shows the two solvers reach
the same accuracy, and that Anderson needs about 30\% fewer iterations but
still costs twice the wall-clock time, because each of its iterations solves a
batched $(m{+}1)\times(m{+}1)$ least-squares system that the plain iteration
does not pay for. At this problem size the cheaper iteration wins.

\section{Problems Encountered}

The backward solve initially stalled at the 50-iteration cap, taking 47
iterations per batch against 17 for plain iteration, despite solving a
\emph{linear} fixed point that Anderson should handle easily. The cause was
the Tikhonov term in Anderson's least-squares step being a fixed absolute
constant. Late in training most examples are classified correctly and
contribute a near-zero gradient, so the residual matrix $G$ becomes tiny, the
regulariser dominates the normal equations, and the computed weights collapse
towards uniform averaging. Scaling the regulariser by $\operatorname{tr}(GG^{\!\top})/m$
fixed it: backward iterations fell from 47 to 11.7 and epoch time from 22.5\,s
to 10.1\,s. This matters beyond the NumPy model, since the same solver is
planned for the convolutional version where the effect would be larger.

\section{Plan for Next Steps}

Progress Report~II (Dec~1) will cover the convolutional DEQ in PyTorch with a
custom backward pass, the truncated Neumann and Jacobian-free gradient
variants~\cite{geng2021training,fung2022jfb}, the unrolled weight-tied
baseline, and the first accuracy, memory and time comparison on CIFAR-10.
Given that Anderson's per-iteration cost dominated here, the comparison will
report wall-clock time alongside iteration counts throughout.

\printrefs
\end{document}
